% title: 2026 春复变函数 (H) 期末考试
% author: 沈维孝
% summary: 从香蕉空间旧站“试卷：复变函数”迁移的完整试卷题目。
% source: https://bananaspace.org/legacy/wiki/用户:Solution/_试卷:_复变函数
\section*{2026 春复变函数 (H) 期末考试（沈维孝）}
\begin{enumerate}
    \item 设$f_{1}, f_{2}$在$\mathbb{D}$上全纯，且$|f_{1}|^{2}+|f_{2}|^{2}\equiv 1$，证明它们为常值函数.
    \item 考虑函数
    $$
    f(z)=\frac{\sin z\cos z}{z^{7}-z^{5}+z^{3}-z}.
    $$
    \begin{itemize}
        \item[(1)] 求$f$的所有零点和极点，并确定其阶数；
        \item[(2)] 计算积分$\displaystyle\int_{|z|=2}\frac{f'(z)}{f(z)}\,dz$.
    \end{itemize}
    \item 设$f$为$\mathbb{C}$上的半纯函数，且满足
    $$
    f\left(\frac{1}{n}\right)=\frac{n^{6}-1}{16n^{4}-1},\qquad n=1,2,\cdots
    $$
    \begin{itemize}
        \item[(1)] 求$f$的解析表达式；
        \item[(2)] 设$f$在$z=\mathrm{i}$处的Taylor展开为$f(z)=\sum a_{n}(z-\mathrm{i})^{n}$，求$\limsup\limits_{n\to\infty}|a_{n}|^{1/n}$.
    \end{itemize}
    \item 设$f:\mathbb{C}\setminus[-1,1]\to\mathbb{C}$为全纯函数，满足
    $$
    f(z)^{2}=z^{2}-1,\qquad f(2)=\sqrt{3}.
    $$
    \begin{itemize}
        \item[(1)] 证明这样的$f$存在；
        \item[(2)] 求$f(-2)$；
        \item[(3)] 求$f$在$\infty$处的留数.
    \end{itemize}
    \item 设$f:\mathbb{D}\to\mathbb{D}$全纯，定义
    $$
    g_{n}(t)=f(r_{n}e^{\mathrm{i} t}),\qquad r_{n}=\frac{n-1}{n},\quad n=1,2,\cdots
    $$
    已知$g_{n}$在$t\in[0,\pi/2]$上一致收敛到$0$，证明$f\equiv 0$.
    \item 考虑半圆盘
    $$
    \Omega=\mathbb{D}\cap\{z\in\mathbb{C}:\Im z>0\}.
    $$
    设$\mathcal{F}$为所有满足$f(0)=\mathrm{i}/2$的全纯函数$f:\mathbb{D}\to\Omega$的集合，求 $\sup_{f\in\mathcal{F}}|f'(0)|$.
\end{enumerate}
